\section{Restoration Tasks and Applications}
\label{sec:tasks}
\subsection{Natural-Video Super-Resolution}
VSR is a particularly clear setting for studying generative ambiguity. A low-resolution frame constrains coarse structure but leaves many high-resolution textures compatible with the measurement. Neighboring frames can reduce this uncertainty through subpixel motion and repeated observation, yet they do not reveal every missing detail. Diffusion-based VSR must decide when to aggregate measured evidence and when to synthesize a plausible completion.

\textbf{Temporalizing image priors.} Upscale-A-Video adds temporal modeling to both denoising and decoding and propagates latents using optical flow~\cite{upscale}. StableVSR couples frames through temporal conditioning and texture guidance, with bidirectional framewise sampling~\cite{stablevsr}. SATeCo uses spatial feature adaptation and temporal feature alignment while retaining pretrained components~\cite{sateco}. MGLD-VSR places a motion constraint directly in the latent sampling path and improves the decoder separately~\cite{mgld}. These methods illustrate different locations for the same broad objective: preventing independently synthesized detail from changing arbitrarily across time.

Later image-prior approaches also address training stability and degradation mismatch. DiffVSR combines progressive learning with interweaved latent transitions~\cite{diffvsr}. The self-supervised ControlNet component in SCST targets degradation-insensitive conditioning, while spatiotemporal Mamba modeling aggregates video information~\cite{scst}. The auxiliary control objective complements paired restoration training by improving the representation of degraded observations.

\textbf{Adapting video priors.} STAR couples a video diffusion prior to low-resolution evidence and uses dynamic frequency control to balance restored structure and detail~\cite{star}. The method is developed with I2VGen-XL/VEnhancer initialization and additionally investigated with CogVideoX for scaling. DC-VSR and DAM-VSR both build on video priors but emphasize different constraints: attention propagation across spatiotemporal partitions in the former, and separation of appearance enhancement from motion control in the latter~\cite{dcvsr,damvsr}. RealisVSR and STCDiT further illustrate Wan-based restoration, including consistency-preserving conditioning and explicit handling of temporal latent compression~\cite{realis,stcdit}.

\textbf{Scaling spatial context.} A model's native training resolution is not necessarily the output resolution needed for restoration. PatchVSR combines local patches with resized global video context so that high-resolution generation is not based solely on isolated tiles~\cite{patch}. This addresses a different problem from merely enlarging the output canvas. Global context helps preserve object-level structure, while local processing controls memory. Seam quality, repeated texture, and consistency of objects crossing patch boundaries remain necessary evaluation targets. Table~\ref{tab:vsr} compares representative designs across prior sources, observation conditioning, temporal mechanisms, and inference regimes.

\begin{table*}[t]
\caption{Representative VSR and generic restoration designs. Backbone labels refer to the cited configurations, not all variants. ``One'' refers to the reported sampling step count; network evaluations and other pipeline costs require separate accounting.}
\label{tab:vsr}\centering\footnotesize
\begin{tabularx}{\textwidth}{@{}p{0.18\textwidth}p{0.19\textwidth}Yp{0.12\textwidth}@{}}\toprule
Method & Prior / initialization & Main observation or temporal mechanism & Inference regime \\\midrule
StableVSR~\cite{stablevsr} & SD $\times4$ image upscaler & Flow-aligned texture guidance and framewise bidirectional sampling & Iterative \\
Upscale-A-Video~\cite{upscale} & SD $\times4$ image upscaler & Temporal U-Net/decoder layers and recurrent latent propagation & Iterative \\
MGLD-VSR~\cite{mgld} & SD 2.1 image prior & Motion-guided latent sampling and temporal-aware decoder & Iterative \\
SATeCo~\cite{sateco} & Stable Diffusion image prior & Spatial adaptation and tubelet temporal alignment & Iterative \\
DiffVSR~\cite{diffvsr} & SD $\times4$ image upscaler & Progressive learning and interweaved latent transition & Iterative \\
STAR~\cite{star} & I2VGen-XL / VEnhancer initialization & Low-resolution information enhancement and dynamic frequency loss & Iterative \\
DC-VSR~\cite{dcvsr} & SVD video prior & Spatial and temporal attention propagation across tiles & Iterative / tiled \\
DAM-VSR~\cite{damvsr} & SVD + reference-image restoration & Separate appearance and motion control & Iterative \\
PatchVSR~\cite{patch} & Pretrained video DiT & Local/global dual-stream conditioning and joint patch modulation & Patchwise \\
DOVE~\cite{dove} & CogVideoX-5B & Latent-then-pixel training for observation-conditioned restoration & One \\
DLoRAL~\cite{dloral} & SD 2.1 image prior & Consistency/detail LoRA separation and cross-frame retrieval & One \\
SeedVR2~\cite{seedvr2} & Trained SeedVR teacher & Adaptive windows, distillation and adversarial post-training & One \\
FlashVSR~\cite{flash} & Wan2.1-1.3B & Streaming distillation, sparse attention and conditional decoder & One / streaming \\
Stream-DiffVSR~\cite{stream} & Distilled SD $\times4$ upscaler & Motion-aligned autoregressive guidance and temporal decoding & Four / frame-causal \\
InfVSR~\cite{infvsr} & Wan2.1-1.3B & Causal DiT, rolling cache and joint visual guidance & One / chunked \\\bottomrule
\end{tabularx}
\end{table*}

\subsection{Arbitrary-Scale and Space-Time Super-Resolution}
Spatial scale and temporal sampling are distinct sources of missing information. AVSR-Diff separates fixed-resolution latent denoising from continuous-coordinate rendering, using flow- and deformable-convolution-based recurrence to align ControlNet residuals and a scale-aware Fourier refinement in the decoder~\cite{avsrdiff}. This design permits flexible output scales without rerunning a resolution-specific diffusion model, but the underlying SD $\times4$ prior still bounds its native generative information: larger outputs require decoder extrapolation. Its 50-step DDPM schedule also places it in a different efficiency regime from contemporary one-step systems. Space-time VSR must additionally infer intermediate observations and motion. OSDEnhancer specializes temporal-coherence and texture-enrichment LoRAs and employs a bidirectional deformable VAE decoder~\cite{osdenhancer}. DiffST combines cross-frame context aggregation for intermediate-frame synthesis with guidance from global video representations~\cite{diffst}. Both adapt video priors for one-step space-time restoration, yet their objectives extend beyond fixed-frame-rate VSR. Spatial magnification and temporal upsampling consequently define separate task dimensions, with intermediate-frame reconstruction relying on motion evidence between the observed frames.

\subsection{Generic, Mixed-Degradation, and Zero-Shot Restoration}
Generic restoration addresses observations for which a single nominal task, such as bicubic super-resolution, is inadequate. Blur, noise, resampling, sharpening, and compression can interact, and their order matters. A model trained on mixtures can amortize restoration across these conditions, but its supported degradation family remains defined by its training and tests.

SeedVR uses a video restoration transformer with windowed spatiotemporal processing to accommodate varying video sizes~\cite{seedvr}. SeedVR2 compresses a stronger restoration process through distillation and adversarial post-training~\cite{seedvr2}. TDM explicitly studies an all-in-one real-world restoration setting with temporal constraints, while Vivid-VR transfers concepts from a text-to-video diffusion transformer~\cite{tdm,vividvr}. These approaches differ in the source of their generative knowledge and in the extent of task-specific learning; ``generic'' does not erase these differences.

The zero-shot route reuses a prior or a pretrained image restoration model and introduces temporal constraints at inference. ZVRD combines temporal attention, consistency guidance, and noise sharing; DiffIR2VR-Zero uses latent warping and token merging~\cite{zvrd,diffir2vr}. DiTVR extends this direction to transformer priors with trajectory-aware attention and flow-consistent sampling~\cite{ditvr}. The journal extension of GDP also treats unified image and video restoration~\cite{gdp}; it should be distinguished from the earlier image-focused GDP paper~\cite{gdp_image}.

The attraction of zero-shot transfer is adaptability without collecting a new paired training set for every task. Its costs can include repeated optimization, strong assumptions about degradation, and reliance on correspondence extracted from corrupted inputs. A zero-shot method is therefore best evaluated together with its operator knowledge, preprocessing, and per-video optimization budget. A frozen model does not imply zero data dependence: the pretrained prior itself embodies a potentially extensive training corpus.

ZVRV~\cite{videodiffassist} uses pretrained video models to assist an image restoration model, rather than asking one backbone to supply both measurement fidelity and temporal coherence. Models with a shared VAE can exchange latents directly, whereas models with different VAEs exchange clean predictions through decoding and re-encoding. A further SVD stage strengthens temporal consistency. The fusion ratio explicitly controls how much the image-restoration trajectory is modified by the video prior. The paper's ``CoT-based'' rule operationally selects among candidate ratios using CLIP-IQA and warping-error ranks; it is a metric-guided best-of-$N$ search. This selection adjusts the contribution of the video prior by rewarding perceptual and temporal proxy scores. Source fidelity requires independent assessment, and candidate evaluation adds inference cost.

\begin{table*}[t]
\caption{Diffusion-based restoration across tasks and application domains.}
\label{tab:domains}\centering\footnotesize
\begin{tabularx}{\textwidth}{@{}p{0.16\textwidth}>{\raggedright\arraybackslash}p{0.22\textwidth}YY@{}}\toprule
Domain / task & Example & Role of generative prior & Task-specific requirement \\\midrule
Deblurring & VD-Diff~\cite{vddiff} & Compact latent prior with wavelet-aware video modeling & Sharp structure must respect exposure and motion \\
Weather & Diffusion TTA; WeatherWeaver~\cite{difftta,weatherweaver} & Conditional diffusion and test-time calibration; video-prior weather control & Preserve scene structure under changing atmospheric effects \\
Low light / events & Wavelet diffusion; EVDiffuser~\cite{lowlight,evdiffuser} & Conditional restoration with frequency or event guidance & Recover weak signal while preserving motion and illumination \\
Underwater & UnDIVE~\cite{undive} & Diffusion-pretrained encoder transferred to enhancement & Physics-based enhancement with a pretrained representation \\
Infrared / turbulence & HATIR~\cite{hatir} & Image diffusion prior with domain-specific guidance & Thermal structure and atmospheric deformation require separate tests \\
Face video & StableBFVR; FLAIR~\cite{stablebfvr,flair} & Generative face prior and temporal conditioning & Identity, expression and motion must accompany naturalness \\
Reference-guided face & IP-FVR~\cite{ipfvr} & Reference identity and identity-specific adaptation & Extra reference/training access changes the comparison setting \\
Codec artifacts & DiQP; PixRevive; DiffCVE~\cite{diqp,pixrevive,diffcve} & Compression-aware process or feature/codec-guided diffusion & Account for quantization and codec side information \\
Inpainting & VipDiff; FFF-VDI; HomoGen~\cite{vipdiff,fffvdi,homogen} & Completion with latent alignment or propagated visual priors & Preserve visible content and reconstruct coherent occlusions \\
HDR video & DiffHDR~\cite{diffhdr} & Video flow prior with radiance inpainting and exposure guidance & Preserve radiance across re-exposure and time \\
Generated video & VEnhancer; CreativeVR~\cite{venhancer,creativevr} & Generative space--time or structure/motion refinement & No unique clean reference; specify permitted content changes \\
Agentic systems & MoA-VR; VQ-Jarvis~\cite{moavr,vqjarvis} & Tool selection can include diffusion operators & Coordinate operators while preserving video content \\\bottomrule
\end{tabularx}
\end{table*}

Table~\ref{tab:domains} relates the role of the generative prior to the requirements of each application.

\subsection{Deblurring and Adverse Weather}
\textbf{Deblurring.} Motion blur integrates scene radiance over an exposure interval. The target sharp frame is therefore linked to exposure timing, camera motion, object motion, and occlusion. A generative prior can supply high-frequency structure, but it must not arbitrarily change the phase of motion or invent crisp edges inconsistent with the integrated observation. VD-Diff combines a wavelet-aware dynamic transformer with a compact diffusion prior~\cite{vddiff}. CPGD-Net adds information from video coding, including motion vectors and coding residuals~\cite{cpgdnet}. Such side information can constrain temporal reconstruction, while also restricting applicability to inputs for which the required bitstream information is available.

\textbf{Weather removal.} Rain streaks, fog, snow, and adherent drops differ in their temporal structure and image formation. A precipitation particle may move independently of the background, while atmospheric scattering changes scene contrast in a depth-dependent way. Treating these effects as generic additive noise can erase useful scene content or leave temporally unstable residuals.

Diffusion test-time adaptation addresses video adverse-weather removal~\cite{difftta}, WeatherWeaver explicitly addresses controllable weather synthesis and removal with video diffusion~\cite{weatherweaver}, and ZSVD explores zero-shot video deraining with video diffusion models~\cite{zsvd}. These approaches establish several routes to weather removal, while open challenges include real mixed weather, spatially varying transmission, night scenes, domain shift, and preservation of small moving objects. The distinction between synthesizing a weather effect and correctly removing it from a real recording should remain explicit.

\subsection{Low-Light and Event-Guided Enhancement}
Low-light video combines weak signal, sensor noise, color uncertainty, and motion blur. Brightening the frames can expose noise, while temporal denoising can smear moving structures. A diffusion prior may recover plausible detail but can also amplify uncertainty about objects that were barely measured. The wavelet-based conditional diffusion approach to low-light video enhancement places low-frequency reconstruction and high-frequency interscale attention in a coordinated pipeline~\cite{lowlight}. EVDiffuser uses event guidance and modality-adaptive processing for consistent video enhancement~\cite{evdiffuser}.

DL-Diff~\cite{dldiff} combines an SD1.5-based spatial prior with a restoration branch and AnimateDiff-style temporal modules, trained progressively on paired frames, normal-light videos, and paired low-/normal-light clips. Although framed around face enhancement, the reported DID and SDSD experiments primarily characterize low-light video quality, leaving identity-specific behavior less fully explored.

Events and RAW measurements provide complementary evidence, but the presence of a temporal sensor does not automatically make a paper a video-restoration study. NEC-Diff investigates event--RAW complementarity in extreme darkness~\cite{necdiff}; claims about long-video consistency require the corresponding sequence-level experiments. Event-guided HDR reconstruction demonstrates diffusion priors for reconstructing HDR images from image/event input~\cite{eventhdr}. We include it as adjacent evidence for conditioning and radiometric reconstruction, rather than treating it as proof of temporally validated HDR video restoration.

The GDP journal extension applies guided diffusion to low-light and LDR video enhancement, using LoLi-phone sequences and videos from the NTIRE 2021 Multi-Frame HDR Challenge~\cite{gdp}. It illustrates the reuse of an image prior across exposure-related video tasks, complementing the video-prior approach described below.

\subsection{High-Dynamic-Range Video Reconstruction}
DiffHDR~\cite{diffhdr} formulates LDR-to-HDR conversion as radiance inpainting, using a Log-Gamma representation to place HDR content within the operating range of an LDR-pretrained video VAE. A VACE/Wan video prior is adapted with LoRA and a rectified flow-matching objective, using training clips synthesized from HDR environment maps. Exposure masks and region-specific text or reference-image conditions guide recovery of clipped highlights and shadows. Radiance behavior across re-exposures and temporal stability in saturated regions are central objectives in this formulation. Moreover, controlled synthesis in a clipped region can be perceptually convincing without recovering its unique original radiance. Reporting controllability and radiometric fidelity as separate properties makes this distinction visible.

HDR also requires a distinction between reconstructing scene radiance, inverse tone mapping, and producing a display-ready enhanced video. A pleasing tone-mapped output is not necessarily radiometrically accurate. Future comparisons should specify transfer functions, exposure, color space, clipping, and display conditions. Temporal brightness stability should be evaluated alongside recovery of genuine illumination changes; a constant exposure curve can be stable yet incorrect.

\subsection{Underwater Enhancement}
Underwater observations combine wavelength-dependent attenuation, backscatter, illumination variation, and often weak texture for motion estimation. Consequently, the color distribution of natural terrestrial videos is not a neutral prior for this domain. Restoring contrast while imposing an attractive terrestrial palette can reduce physical fidelity.

UnDIVE is informative because it uses generative pretraining to obtain a transferable representation and then employs a physics-based enhancement formulation and temporal consistency learning~\cite{undive}. Its diffusion-derived encoder should not be described as an ordinary reverse-diffusion video sampler. The example broadens the survey's definition of a useful diffusion prior while reinforcing the need to report the actual inference graph.

Evaluation is difficult when clean underwater references are unavailable. Visibility, color plausibility, and temporal consistency answer different questions, and a no-reference underwater score cannot establish recovery of true scene colors. A stronger protocol would combine controlled capture where possible, task-level measurements such as tracking or inspection accuracy, and human assessment with explicit instructions about color and visibility. Changes in water type, lighting, depth, and camera response should be treated as domain shifts rather than hidden within one mean quality score.

\subsection{Infrared and Atmospheric Turbulence}
Thermal infrared video and turbulence mitigation require more than transferring an RGB texture prior. Thermal intensity relates to emission, transmission, sensor calibration, and scene material, while turbulence introduces time-varying geometric distortions and blur. A restored sequence can appear sharp while corrupting temperature-related structures or converting random atmospheric deformation into apparent object motion.

HATIR studies turbulent infrared VSR with heat-aware guidance, motion-related modeling, and turbulence-aware reconstruction~\cite{hatir}. Its domain-specific design connects an image prior to the thermal signal and the deformation process. Generalization across sensors and turbulence strengths remains challenging, particularly when atmospheric distortion must be separated from genuine object motion. Structural recovery and radiometric preservation are distinct objectives in this setting.

Turbulence also exposes a limitation of generic temporal smoothing: the correct output should remove atmospheric jitter while retaining scene motion. A method that reduces all motion may score well under a naive consistency metric. Ground-truth trajectories, static and dynamic scene subsets, and confidence-aware alignment provide more discriminating tests.

\subsection{Face Video Restoration}
Face restoration introduces identity and expression as explicit fidelity targets. Fine texture can be plausible without belonging to the recorded person, and a temporally consistent face can still have incorrect geometry. StableBFVR incorporates a temporally consistent generative prior and degradation-aware prompting~\cite{stablebfvr}. FLAIR adapts an image diffusion prior through recurrent refinement and multiscale temporal attention, addressing several face-video restoration tasks~\cite{flair}. Reference-guided IP-FVR uses additional identity information and identity-specific LoRA adaptation, while DVFace explores spatiotemporal priors in one-step video face restoration~\cite{ipfvr,dvface}.

Reference-based and reference-free methods should not share an undifferentiated leaderboard. A high-quality portrait is additional information that can materially reduce uncertainty. Evaluation should report reference identity, acquisition conditions, and whether the reference could be available in the intended use case. Identity embeddings should be supplemented by expression, pose, and temporal inspection, since one feature score may tolerate errors important to human viewers.

For archival, forensic, or biometric applications, synthesized detail must be distinguished from recovered evidence. A system can expose uncertainty or retain the original video alongside the restoration. The scientific objective is to improve visible content without presenting prior-driven facial features as verified observations.

\subsection{Compression Artifacts and Bitstream Corruption}
Compression artifacts are structured by the codec, quantization parameters, prediction modes, and reference frames. They differ from independent Gaussian noise. DiQP uses a QP-aware transformer--diffusion formulation for high-resolution codec-degraded video; its corruption process is closely tied to compression rather than standard Gaussian noising~\cite{diqp}. Metadata-guided diffusion has also been studied for recovering bitstream-corrupted video~\cite{bitstream}.

The restoration problem depends on codec version, chroma sampling, quantization, and available side information. In post-processing, the transmitted representation is already fixed; a generative codec additionally learns which information to transmit. This difference changes both the observation model and the quality--rate tradeoff.

Compressed-video enhancement exploits both inter-frame redundancy and the structure introduced by coding. PixRevive~\cite{pixrevive} aligns neighboring compressed frames, extracts wavelet-enriched latent features, and uses these features to condition a task-built diffusion denoiser with group-wise feature fusion. Here, ``latent features'' denotes the ELPNet conditioning features, rather than a latent space supplied by a pretrained VAE. DiffCVE~\cite{diffcve} instead adapts SD2.1 and explicitly uses codec residuals, motion vectors, and QP-dependent semantic conditions. Temporal processing appears in both the denoiser and VAE decoder, while coding-prior features are fused with QP-dependent weights. Its current version reports both a 50-step model and a separately trained fixed-timestep one-step model. The one-step variant is learned with fixed-timestep paired training, illustrating a route to direct restoration distinct from teacher--student distillation. Access to codec-side information further separates this setting from enhancement based only on decoded RGB.

Relic et al.~\cite{stdpevc} move the diffusion prior inside the codec itself. Their SVD-based decoder reconstructs a GOP from transmitted keyframes and compressed sparse bidirectional flow/mask information, using long-context interpolation to synthesize detail and motion. This is an adjacent generative-coding problem: the encoder selects which information to transmit, whereas post-decoding enhancement accepts a fixed bitstream produced elsewhere. Rate, all side-information bits, coding delay, and reconstruction quality must be compared jointly. Its experiments emphasize very low-rate compression, where a strong generative decoder can replace a substantial portion of transmitted appearance detail with prior information.

Temporal behavior matters particularly near intra-refresh boundaries, scene cuts, packet loss, and changes in quantization. Restored texture that is sharper in selected frames but unstable across a group of pictures can be less useful than a modest, consistent correction. For pipelines that decode and re-encode a stream, these stages contribute to end-to-end streaming cost.

\subsection{Film Restoration}
HaineiFRDM~\cite{haineifrdm} addresses a different source of ambiguity: a local patch in an old film may not reveal whether a structure is a defect or part of a fast-moving object. It augments an SD1.5/AnimateDiff initialization with patch-position information, global frame and prompt context, frequency-based texture features, and patch-consistent inference. The combination of local texture and global context helps distinguish film defects from moving scene content while maintaining consistency across patches.

\subsection{Inpainting and Generated-Video Enhancement}
\textbf{Inpainting.} Video inpainting treats a masked region as unobserved. The correct target may be ambiguous, especially when a region is never visible in any frame. VipDiff uses training-free diffusion with temporal constraints, AVID addresses videos beyond a fixed native length, and DiffuEraser combines diffusion with a coherent initialization and expanded temporal support~\cite{vipdiff,avid,diffueraser}. FloED uses flow guidance and efficiency mechanisms for coherent completion~\cite{floed}. These methods illustrate why diversity and consistency can coexist: multiple completions may be acceptable if each follows the visible geometry and motion.

FFF-VDI~\cite{fffvdi} uses future-frame latent information to fill the first-frame conditioning signal of an SVD-based image-to-video model, then refines noisy latent alignment with a learned deformable module. Temporal transformer fine-tuning adapts the pretrained motion representation to masked observations. HomoGen~\cite{homogen} makes a different choice about the reliability of propagated evidence. Bidirectional homography propagation provides contextual pixels, while a content-adaptive diffusion controller treats these pixels as fallible priors rather than final reconstructions. Learned scaling and injection let the denoiser use their semantic structure while correcting warping artifacts. These designs illustrate why propagation accuracy and generative capacity must be evaluated together: retaining more visible evidence is useful only if the reconstruction model can recognize and correct propagation errors.

Infusion~\cite{infusion} targets the complementary setting where the observed video itself contains sufficient repeated structure. It trains a small pixel-space 3D diffusion network on that single video, excluding the actual test-mask region from supervision. Interval training alternates optimization and inference over portions of the noise schedule. This approach trades a large external prior for per-video computation and an assumption of internal self-similarity; it is especially suited to recurring dynamic textures, while novel semantic content hidden throughout the clip remains difficult. EraserDiT~\cite{eraserdit} instead transfers an LTX-Video transformer and fine-tunes it with masked videos and prompts. Its Circular Position-Shift changes the temporal windows seen during denoising to reduce inconsistencies in long sequences. HomoGen's overlapping GOP references and EraserDiT's shifting windows both extend finite-context inference by carrying information across local processing boundaries.

Masked-region quality and unmasked-region preservation should be reported separately. A method can improve the apparent naturalness of a filled region by unnecessarily changing visible content. Likewise, mask-aware flow estimation and disocclusion behavior deserve explicit analysis rather than relying only on whole-frame averages.

\textbf{AIGC enhancement.} Generated videos may contain local texture defects, inconsistent object structure, implausible motion, or semantic errors that cannot be modeled as simple downsampling. VEnhancer performs generative space--time enhancement~\cite{venhancer}; Noise Calibration constrains content changes during enhancement~\cite{ncsdedit}; CreativeVR and GenVE address structural, motion, semantic, or texture aspects of generative-video refinement~\cite{creativevr,genve}. Here ``fidelity'' includes preservation of intended content, not necessarily recovery of a historically observed scene.

Without a unique clean reference, enhancement is constrained by attributes such as identity, object count, camera path, action order, and prompt-relevant semantics. Smoother motion is not always a correction if it changes the requested action. User preference, content preservation, and physical plausibility should therefore be measured as separate dimensions.

Generated-video enhancement changes the degradation model as well as the input distribution. In a cascaded generator, the low-resolution input can remain in latent space, and its errors need not resemble blur, noise, and compression applied to a real video. SimpleGVR addresses this interface through a latent upsampler and flow-based and model-guided degradations that better match upstream text-to-video outputs~\cite{simplegvr}. A detail-aware training timestep sampler emphasizes difficult reconstruction stages, while interleaving temporal units supports long-sequence processing. Structural defects introduce a further distinction between detail recovery and geometric correction. DTG-Restore refines a video using a cleaner diffusion-time anchor and a decaying time offset, then optionally applies an off-the-shelf restoration module~\cite{dtgrestore}. This guidance couples estimates at different diffusion times while using the distorted clip as conditioning for geometric refinement. GenWarp480 collects generated clips with characteristic geometric artifacts, providing a complementary setting to conventional synthetic degradation benchmarks. Evaluation in this setting should separate preservation of intended content, correction of implausible geometry, and the perceptual appeal of newly synthesized detail.
